\chapter{Descriptions textuelles des cas d'utilisation}\label{chap:cas}

Les messages cités sont ceux affichés par l'application. Sauf mention contraire, chaque refus réaffiche la
page avec un message d'erreur et ne modifie aucune donnée.

\section{S'inscrire (BF01)}

\fichecas{Visiteur}
  {Le visiteur n'est pas connecté.}
  {Un compte CLIENT actif est créé, mot de passe haché.}
  {RG01, RG02, BNF01}

\scenario{Scénario nominal}
\begin{enumerate}
  \item Le visiteur ouvre la page « Inscription » (\code{/inscription}).
  \item Il saisit nom, prénom, e-mail, mot de passe, confirmation, et facultativement adresse et téléphone.
  \item Le système vérifie les champs obligatoires, le format de l'e-mail, la longueur minimale du mot de passe
        (6 caractères), la longueur des champs et l'égalité des deux mots de passe.
  \item Le système vérifie que l'e-mail (enregistré en minuscules, sans espaces) n'est pas déjà utilisé (RG01).
  \item Le système crée le compte avec le rôle CLIENT (RG02) et le mot de passe haché par BCrypt (BNF01).
  \item Le visiteur est redirigé vers la page de connexion : « Inscription réussie. Vous pouvez vous connecter. »
\end{enumerate}

\scenario{Scénarios alternatifs et d'erreur}
\begin{itemize}
  \item 3a. Champ obligatoire vide : « Le nom, le prénom, l'e-mail et le mot de passe sont obligatoires. »
  \item 3b. E-mail mal formé : « L'adresse e-mail n'est pas valide. »
  \item 3c. Mot de passe trop court : « Le mot de passe doit contenir au moins 6 caractères. »
  \item 3d. Confirmation différente : « Les deux mots de passe ne correspondent pas. »
  \item 3e. Champ trop long : « Un des champs saisis est trop long. »
  \item 4a. E-mail déjà utilisé : « Cette adresse e-mail est déjà utilisée. »
  \item Dans tous ces cas, le formulaire est réaffiché avec les valeurs saisies (sauf le mot de passe).
\end{itemize}

\section{S'authentifier et se déconnecter (BF02, BF03)}

\fichecas{Visiteur (connexion) ; Client, Administrateur (déconnexion)}
  {Le visiteur n'est pas connecté.}
  {Une nouvelle session contient l'utilisateur connecté ; à la déconnexion, la session est invalidée.}
  {BNF01, BNF04}

\scenario{Scénario nominal (connexion)}
\begin{enumerate}
  \item Le visiteur ouvre « Connexion » (\code{/connexion}) et saisit son e-mail et son mot de passe.
  \item Le système recherche l'utilisateur par e-mail et compare le mot de passe au hachage BCrypt.
  \item Le système vérifie que le compte est actif.
  \item L'ancienne session est invalidée, une nouvelle session est créée avec l'utilisateur (protection contre
        la fixation de session).
  \item Un client est redirigé vers le catalogue, un administrateur vers le tableau de bord.
\end{enumerate}

\scenario{Scénario nominal (déconnexion)}
\begin{enumerate}
  \item L'utilisateur connecté clique sur « Déconnexion » (formulaire POST vers \code{/deconnexion}).
  \item La session est invalidée (le panier est donc vidé) et l'utilisateur est redirigé vers la page de
        connexion : « Vous êtes déconnecté. »
\end{enumerate}

\scenario{Scénarios alternatifs et d'erreur}
\begin{itemize}
  \item 1a. Champ vide : « Veuillez saisir votre e-mail et votre mot de passe. »
  \item 2a. E-mail inconnu ou mot de passe faux : « E-mail ou mot de passe incorrect. » (même message dans les deux
        cas, pour ne pas révéler l'existence d'un compte).
  \item 3a. Compte désactivé par l'administrateur : « Ce compte est désactivé. »
  \item Session inactive pendant 30 minutes : elle expire et l'utilisateur doit se reconnecter (BNF04).
\end{itemize}

\section{Consulter et rechercher le catalogue (BF05 à BF08)}

\fichecas{Visiteur (et donc Client)}
  {Aucune.}
  {Aucune modification des données.}
  {RG06, RG15, BNF10}

\scenario{Scénario nominal}
\begin{enumerate}
  \item Le visiteur ouvre le catalogue (\code{/catalogue}, page d'accueil de l'application).
  \item Le système affiche les produits actifs (RG15), les plus récents d'abord, par pages de 10 (BNF10), avec
        image, nom, catégorie et prix. Un produit à stock 0 reste affiché avec le badge « Rupture » (RG06).
  \item Le visiteur choisit une catégorie (BF06) et/ou saisit un mot-clé (BF07) : le système filtre sur la
        catégorie et cherche le mot-clé dans le nom et la description, sans tenir compte de la casse. Les filtres
        sont conservés dans les liens de pagination.
  \item Le visiteur ouvre la fiche d'un produit (\code{/produit?id=N}, BF08) : image, prix, stock, description,
        catégorie.
\end{enumerate}

\scenario{Scénarios alternatifs et d'erreur}
\begin{itemize}
  \item 2a. Aucun produit ne correspond : la liste est vide avec un message.
  \item 2b. Numéro de page invalide ou hors limites : le système affiche la page existante la plus proche.
  \item 4a. Produit inconnu, désactivé ou identifiant invalide : retour au catalogue avec « Produit introuvable. »
  \item 4b. Sur la fiche, un visiteur voit « Connectez-vous pour commander » ; un client voit le formulaire
        « Ajouter au panier » si le produit n'est pas en rupture.
\end{itemize}

\section{Gérer le panier (BF09 à BF12)}

\fichecas{Client}
  {Le client est connecté (RG08, contrôlé par \code{AuthFilter}).}
  {Le panier de la session est à jour ; rien n'est enregistré en base.}
  {RG06, RG07, RG08, RG15}

\scenario{Scénario nominal}
\begin{enumerate}
  \item Sur la fiche d'un produit, le client choisit une quantité et clique sur « Ajouter au panier » (BF09).
  \item Le système recharge le produit depuis la base, vérifie qu'il est actif, que la quantité est un entier
        supérieur ou égal à 1, que le stock n'est pas nul (RG06) et que la quantité déjà au panier plus la
        quantité ajoutée ne dépasse pas le stock (RG07).
  \item Le produit est ajouté au panier (ou sa quantité est augmentée) : « Produit ajouté au panier. »
  \item Le client ouvre « Panier » (\code{/client/panier}, BF12) : lignes, sous-totaux et montant total.
  \item Il modifie la quantité d'une ligne (BF10, mêmes contrôles RG06 et RG07), retire une ligne ou vide le
        panier (BF11).
\end{enumerate}

\scenario{Scénarios alternatifs et d'erreur}
\begin{itemize}
  \item 1a. Visiteur non connecté : redirection vers \code{/connexion}. Administrateur : redirection vers le tableau de
        bord (seul un client a un panier).
  \item 2a. Produit inconnu ou désactivé : « Produit introuvable. »
  \item 2b. Quantité vide, nulle, négative ou non numérique : « Quantité invalide : indiquez un nombre entier
        supérieur ou égal à 1. »
  \item 2c. Stock nul : « Produit en rupture de stock. »
  \item 2d. Quantité cumulée supérieure au stock : « Quantité demandée supérieure au stock disponible (X en stock,
        Y déjà dans votre panier). »
  \item Après un refus à l'ajout, le client revient sur la fiche du produit avec le message.
\end{itemize}

\section{Passer une commande (BF13)}

\fichecas{Client}
  {Le client est connecté et son panier n'est pas vide.}
  {Une commande EN\_ATTENTE et ses lignes sont enregistrées, les stocks sont décrémentés, le panier est vidé.}
  {RG09, RG10, RG11, RG12, RG15, BNF05}

\scenario{Scénario nominal}
\begin{enumerate}
  \item Depuis le panier, le client clique sur « Passer la commande ».
  \item Le système affiche la page de validation (\code{/client/commandes?action=valider}) : récapitulatif du panier
        et adresse de livraison préremplie avec l'adresse du profil.
  \item Le client vérifie ou modifie l'adresse et clique sur « Confirmer la commande ».
  \item Dans une seule transaction (BNF05), le système recharge le compte du client, puis chaque produit du
        panier avec un verrou en écriture ; il vérifie que le produit est actif et que le stock suffit.
  \item Pour chaque ligne, le stock est décrémenté (RG12) et une ligne de commande est créée avec le prix actuel
        du produit (RG10).
  \item Le montant total est calculé comme la somme quantité × prix unitaire (RG11) ; la commande est enregistrée
        avec ses lignes, au statut EN\_ATTENTE.
  \item Le panier est vidé et le client est redirigé vers « Mes commandes » : « Votre commande n° X a bien été
        enregistrée. Montant total : … DH. »
\end{enumerate}

\scenario{Scénarios alternatifs et d'erreur}
\begin{itemize}
  \item 3a. Panier vide (RG09) : « Votre panier est vide. »
  \item 3b. Adresse vide ou trop longue : « L'adresse de livraison est obligatoire. » / « L'adresse de livraison est
        trop longue (255 caractères maximum). »
  \item 4a. Compte désactivé depuis la connexion : « Votre compte est désactivé : vous ne pouvez plus commander. »
  \item 4b. Produit désactivé ou supprimé depuis l'ajout au panier (RG15) : « Le produit « nom » n'est plus
        disponible : retirez-le de votre panier. »
  \item 4c. Stock insuffisant (RG12) : « Stock insuffisant pour « nom » (X en stock, Y demandé(s)). »
  \item Dans les cas 4a à 4c, la transaction est annulée (rollback) : aucune commande, aucune ligne, aucun stock
        modifié. La page de validation est réaffichée avec le message et l'adresse saisie ; le panier est conservé.
\end{itemize}

\section{Consulter ses commandes et annuler une commande (BF14, BF15, BF16)}

\fichecas{Client}
  {Le client est connecté.}
  {Consultation : aucune modification. Annulation : commande ANNULEE, quantités remises en stock.}
  {RG13, RG14, RG16, BNF10}

\scenario{Scénario nominal}
\begin{enumerate}
  \item Le client ouvre « Mes commandes » (\code{/client/commandes}, BF14) : ses commandes uniquement (RG16), les plus
        récentes d'abord, par pages de 10, avec numéro, date, montant et statut.
  \item Il ouvre le détail d'une commande (\code{/client/commandes?id=N}, BF15) : lignes, prix unitaires enregistrés,
        sous-totaux, total, adresse de livraison et statut.
  \item Si la commande est « En attente », il clique sur « Annuler la commande » et confirme (BF16).
  \item Le système verrouille la commande, vérifie qu'elle appartient au client et qu'elle est EN\_ATTENTE, remet
        chaque quantité en stock (RG14) et passe la commande à ANNULEE, dans une même transaction.
  \item Le détail est réaffiché : « Votre commande n° X a été annulée. »
\end{enumerate}

\scenario{Scénarios alternatifs et d'erreur}
\begin{itemize}
  \item 2a. Commande d'un autre client, identifiant inconnu ou invalide : retour à la liste avec « Commande
        introuvable. » (le système ne révèle pas l'existence de la commande, RG16).
  \item 3a. Commande dans un autre statut : le bouton n'est pas affiché ; une requête forgée est refusée avec
        « Seule une commande « En attente » peut être annulée. » (une commande « Validée » ne peut être annulée
        que par l'administrateur).
  \item 4a. Deux demandes d'annulation simultanées : le verrou les traite l'une après l'autre ; la seconde est
        refusée et le stock n'est réintégré qu'une fois.
\end{itemize}

\section{Consulter et modifier son profil (BF04)}

\fichecas{Client}
  {Le client est connecté.}
  {Les informations (ou le mot de passe) sont mises à jour ; le rôle et l'état du compte ne changent pas.}
  {RG01, BNF01}

\scenario{Scénario nominal}
\begin{enumerate}
  \item Le client ouvre « Mon profil » (\code{/client/profil}).
  \item Il modifie nom, prénom, e-mail, adresse ou téléphone et enregistre.
  \item Le système contrôle les champs et vérifie que le nouvel e-mail n'appartient à aucun autre compte (RG01),
        puis enregistre ; l'utilisateur en session est remis à jour : « Votre profil a été mis à jour. »
  \item Dans le second formulaire, il saisit son mot de passe actuel, le nouveau et la confirmation ; le système
        vérifie le mot de passe actuel puis enregistre le nouveau haché (BNF01) : « Votre mot de passe a été
        modifié. »
\end{enumerate}

\scenario{Scénarios alternatifs et d'erreur}
\begin{itemize}
  \item 3a. Champ obligatoire vide, e-mail invalide ou déjà utilisé, champ trop long : message et formulaire
        réaffiché.
  \item 4a. Mot de passe actuel incorrect : « Le mot de passe actuel est incorrect. »
  \item 4b. Confirmation différente ou nouveau mot de passe de moins de 6 caractères : message d'erreur.
\end{itemize}

\section{Gérer les catégories, les produits et le stock (BF17, BF18, BF19)}

\fichecas{Administrateur}
  {L'administrateur est connecté (\code{AdminFilter}).}
  {Le catalogue est mis à jour.}
  {RG03, RG04, RG05, RG06, RG15}

\scenario{Scénario nominal (catégories, \code{/admin/categories})}
\begin{enumerate}
  \item L'administrateur consulte la liste des catégories.
  \item Il crée ou modifie une catégorie (nom obligatoire et unique, description facultative) : « Catégorie
        enregistrée. »
  \item Il supprime une catégorie vide : « Catégorie supprimée. »
\end{enumerate}

\scenario{Scénario nominal (produits, \code{/admin/produits})}
\begin{enumerate}
  \item L'administrateur consulte la liste paginée de tous les produits, actifs ou non (les inactifs sont grisés).
  \item Il crée ou modifie un produit : nom, description, prix, stock, URL de l'image, catégorie, case « Actif ».
  \item Le système vérifie la catégorie (RG03), le prix strictement positif après arrondi au centime et le stock
        entier positif ou nul (RG05) : « Produit enregistré. »
  \item Il met à jour le stock directement depuis la liste (BF19) : « Stock mis à jour. »
  \item Il supprime un produit : s'il n'a jamais été commandé, il est supprimé ; sinon il est désactivé (RG15) et
        disparaît du catalogue public.
\end{enumerate}

\scenario{Scénarios alternatifs et d'erreur}
\begin{itemize}
  \item Catégorie : nom vide ou déjà pris (« Une catégorie porte déjà ce nom. ») ; suppression d'une catégorie qui
        contient des produits, même inactifs (RG04) : « Cette catégorie contient des produits : elle ne peut pas
        être supprimée. »
  \item Produit : catégorie absente (« Veuillez choisir une catégorie. »), prix nul, négatif ou non numérique
        (« Le prix doit être strictement positif. »), stock négatif (« Le stock doit être un entier supérieur ou
        égal à 0. »), champ trop long.
  \item Un produit désactivé par RG15 peut être réactivé en cochant « Actif ».
  \item Identifiant inconnu ou invalide : « Produit introuvable. » / « Catégorie introuvable. »
\end{itemize}

\section{Gérer les commandes et changer leur statut (BF20, BF21)}

\fichecas{Administrateur}
  {L'administrateur est connecté.}
  {Le statut de la commande est modifié ; en cas d'annulation, les quantités sont remises en stock.}
  {RG13, RG14}

\scenario{Scénario nominal}
\begin{enumerate}
  \item L'administrateur ouvre « Commandes » (\code{/admin/commandes}) : toutes les commandes avec leur client, par
        pages de 10, plus récentes d'abord ; il peut filtrer par statut (BF20).
  \item Il ouvre le détail d'une commande ou reste sur la liste.
  \item Il choisit le nouveau statut parmi ceux autorisés depuis le statut actuel (RG13) : EN\_ATTENTE → VALIDEE
        ou ANNULEE ; VALIDEE → EXPEDIEE ou ANNULEE ; EXPEDIEE → LIVREE.
  \item Le système verrouille la commande, contrôle la transition et enregistre : « La commande n° X est
        maintenant « … ». »
  \item Pour une annulation (confirmation demandée), chaque quantité est remise en stock dans la même transaction
        (RG14) : « … Le stock a été réintégré. »
\end{enumerate}

\scenario{Scénarios alternatifs et d'erreur}
\begin{itemize}
  \item 3a. Commande LIVREE ou ANNULEE : statut définitif, aucun choix proposé.
  \item 4a. Transition interdite (requête forgée : même statut, retour en arrière, saut d'étape) : « Changement de
        statut non autorisé : une commande « … » ne peut pas passer à « … ». » ; statut inconnu : « Statut de
        commande inconnu. »
  \item 4b. Commande inexistante : « Commande introuvable. »
  \item Un client qui accède à \code{/admin/*} reçoit une page d'erreur 403 ; un visiteur est redirigé vers la
        connexion.
\end{itemize}

\section{Gérer les utilisateurs (BF22)}

\fichecas{Administrateur}
  {L'administrateur est connecté.}
  {Le compte est activé ou désactivé.}
  {BNF02 (aucune règle de gestion propre à ce cas)}

\scenario{Scénario nominal}
\begin{enumerate}
  \item L'administrateur ouvre « Utilisateurs » (\code{/admin/utilisateurs}) : liste paginée (nom, e-mail, téléphone,
        rôle, date d'inscription, état), plus récents d'abord.
  \item Il clique sur « Désactiver » (avec confirmation) ou « Activer » pour un compte client.
  \item Le système enregistre : un compte désactivé ne peut plus se connecter ni passer de commande.
\end{enumerate}

\scenario{Scénarios alternatifs et d'erreur}
\begin{itemize}
  \item 2a. Compte administrateur : « Un compte administrateur ne peut pas être désactivé. » (ce qui empêche aussi
        l'administrateur de perdre son propre accès).
  \item 2b. Identifiant inconnu : « Utilisateur introuvable. »
  \item 3a. Un client déjà connecté au moment de sa désactivation garde sa session jusqu'à sa déconnexion ou son
        expiration, mais ne peut plus commander.
\end{itemize}

\section{Consulter le tableau de bord (BF23)}

\fichecas{Administrateur}
  {L'administrateur est connecté (il arrive sur cette page après la connexion).}
  {Aucune modification des données.}
  {RG06, RG15}

\scenario{Scénario nominal}
\begin{enumerate}
  \item L'administrateur ouvre « Tableau de bord » (\code{/admin/tableau-de-bord}).
  \item Le système affiche le nombre total de commandes et le nombre par statut (chaque statut mène à la liste
        filtrée), le chiffre d'affaires (somme des montants des commandes non annulées, commandes en attente
        comprises) et la liste des produits actifs dont le stock est à 0, avec un lien pour les réapprovisionner.
\end{enumerate}

\scenario{Scénarios alternatifs}
\begin{itemize}
  \item 2a. Base vide : 0 commande, « 0,00 DH », « Aucun produit en rupture. »
\end{itemize}
